\subsection{Three Defenses, and Measuring the Slope Not the Level}
\label{sec:9.3.5}
The Three Rs closed on the real risk this whole structure faces: institutions that defend the classification they started with rather than revisit it when evidence changes. A review board, a legal-status framework, and an enforcement mechanism tied to funding eligibility solve the design problem. None of it survives contact with an institution that has stronger incentives to avoid a finding than to make one.

Transparency is the first defense, and it works by making avoidance visible rather than by appealing to good faith. A research program that publishes its protocols, its review-board submissions, and its outcomes — including the ones that didn't work — gives outside researchers, journalists, and future review boards something to check a claim of compliance against. A program that publishes only its successes is not transparent in any sense that matters here; selective disclosure is a more sophisticated way of avoiding scrutiny, not an alternative to it.

The second defense is the one section~\ref{sec:9.3.4}'s enforcement mechanism doesn't cover by itself: protection for the people inside an institution who notice a problem the institution would rather not have found. A funding-eligibility penalty deters an institution from bypassing review, but it does nothing for the individual researcher who reports that a colleague did anyway — and an institution facing that penalty has every incentive to make the reporter's position uncomfortable rather than fix what was reported. Whistleblower protections, and a reporting channel that doesn't route through the people being reported on, are what make that mechanism something a researcher can actually use.

\runin{What a procedure records, and what it drops}
The reporter's position gets uncomfortable before anything formal happens to them, and the incentive on the institution does not fully explain why. Professional culture treats affect as the opposite of judgment — the thing a serious decision-making process succeeds by excluding. Excluding affect does not produce neutrality. It merely leaves the interests that are already written down to govern: revenue, throughput, legal exposure, schedule, institutional survival. Those are the interests a procedure was built to record. Humiliation, dependency, fear of a manager, and the sense that something done to a person was degrading are not less real for having no field on the form — quite the contrary: they are signals through which a harm becomes visible as a harm rather than as a cost.

Treat them as defeasible moral signals — presumptive evidence of harm, real enough to require an answer and defeasible enough to be overruled — and the position is workable. Any such signal can be wrong, and testing it is what deliberation, evidence, and a review process are for. An institution that declines to hear them at all has not achieved impartiality. It has adopted a selection rule that admits whatever it already measures, which is the second of the four structural features section~\ref{sec:2.1.2} uses to define fascism: the measure substituted for the thing measured, so that what falls outside the number is not refuted but unnamed. An institution that shames a reporter does not need to discipline them, because shame construes the self rather than the act as wrong and produces withdrawal.

The third defense (of the oversight machinery against the institution it is supposed to bind) is: standing infrastructure rather than a one-time policy. The Partnership on AI, a voluntary body whose own design limits what it can oblige anyone to do \autocite{partnershiponai2016}, is one example of the kind of body that keeps a set of standards live, revisited as the underlying technology changes. A culture that treats its ethical commitments as a completed project is the same failure mode warned about earlier, applied to institutions instead of individual classifications: having gotten it right once is not the same thing as staying right.

\runin{Measure the slope, not the level}
An institution can keep all three defenses and metabolize all three. The channel exists, the protection is written down, the standing body meets. Dissent is permitted to appear, the institution points at the appearance as proof it listens, and nothing the dissent was meant to move moves — recuperation, the first of the four structural features in that definition of fascism.

There is no threshold that detects this. Any real institution has some overlap between the authority that sets a standard and the authority that grants exceptions to it, and registering an objection is always at least slightly more costly than staying quiet. Asking whether an institution has crossed into capture asks for a line nobody can draw.

Is the cost of registering dissent rising, year over year, while the underlying risk stays flat? Is the required justification getting longer? Are there more approvers than there were? Are the roles that gate an exception migrating toward the roles that set the standard being excepted from? None of those is damning in a single year, and all of them are measurable across several.

The reason to prefer the slope is that it is harder to fake. An institution can present a level — we have a whistleblower channel, we have an ethics board — and the presentation costs nothing. A trend in the cost of using the channel is composed of artifacts the institution generated itself for other reasons, and cannot be curated after the fact without falsifying a record. Where the record is public, the slope is not even self-reported: anyone outside can measure whether dissent has been getting more expensive.

An instrument for detecting recuperation is an unusually easy thing to recuperate. The first caution is that whoever measures the slope must not also gate the exceptions; a review that acquires the power to block becomes the thing it was built to watch. The second is that the measurement has to be refusable — an audit that observes is a different object from a committee that approves, and the same hundred words of written justification are dissent when the dissenter prices them and tribute when the gatekeeper does. The text of the document cannot tell those apart. Only the refusal rights can.
